% Separate checked post-CP19 source. No geometric object atlas is asserted.
\subsection{A weighted component bound for a prospective normal-layer atlas}

\begin{lemma}[Proved: weighted component bound with arbitrary characteristic]
\label{D:weighted-affine-component-bound}
Let $k$ be algebraically closed, let $x_1,\ldots,x_N$ have
positive integral weights $w_1,\ldots,w_N$, and let
$D\geq\max_i w_i$ be an integer. If a closed subscheme
$V\subset\mathbb A^N_k$ is defined by any finite number of
polynomials of weighted degree at most $D$, then
\[
       \#\operatorname{Irr}(V)
          \leq (N+1)\frac{D^N}{\prod_{i=1}^N w_i}.
                                                        \tag{D.WB.1}
\]
The same bound applies to an open subscheme of $V$ and to the
number of nonempty clopen parts in any partition of that open
subscheme. The weights may be divisible by $\operatorname{char}k$.
\end{lemma}

\begin{proof}
We use the ordinary isolated affine intersection bound, including
local scheme lengths, in \citet[Section IV.4.2, equations
(38)--(39), Proposition IV.28, and Corollary VIII.3 following
Theorem VIII.2, pp.~70--71 and 143--144]{Mondal2021WeightedD}.
For $N$ polynomials in $N$ variables of ordinary degrees at
most $D$, the sum of the lengths at their isolated zeroes is
at most $D^N$, even when other components have positive
dimension. The source works over every algebraically closed
field. Its weighted Theorem VIII.2 also directly supplies
the next square-system bound; we record the finite-flat
comparison to make the inseparable-weight scope explicit.

Apply it now to a square weighted system by the finite flat
coordinate-powering morphism
\[
 \mathbb A^N_y\longrightarrow\mathbb A^N_x,
                   \qquad x_i=y_i^{w_i}.                 \tag{D.WB.2}
\]
The monomials $\prod y_i^{a_i}$ with $0\leq a_i<w_i$ form
a basis over $k[x_1,\ldots,x_N]$, so its rank is
$W=\prod_i w_i$, without a separability hypothesis.
A polynomial of weighted degree at most $D$ pulls back to
one of ordinary degree at most $D$. Above an isolated local
zero scheme, the sum of the lengths of the inverse-image local
schemes is exactly $W$ times its length. Thus the number of
isolated geometric points of any square weighted system is
at most $D^N/W$. This includes points with zero coordinates
and inseparable coordinate powers.

We reduce component counting to this square isolated case.
Fix a dimension $d$. Let $I$ be the ideal generated by the
given polynomials. For each $d$-dimensional irreducible
component with generic prime $P$, the regular local ring
$k[x]_P$ has dimension $N-d$, and $I_P$ is primary for its
maximal ideal. Choose $N-d$ constant $k$-linear combinations
of the given generators that form a system of parameters
simultaneously at all these finitely many primes. This
choice is possible by successive prime avoidance: at each
step the bad coefficient choices form finitely many proper
linear conditions over the residue fields, and their
restrictions to the infinite field $k$ do not cover the
coefficient space. Equivalently one can make the choices
over independent generic constants and then specialize in
the resulting nonempty open conditions.

Each selected $d$-dimensional component is consequently an
actual component of the zero scheme of these $N-d$ equations.
Intersect with $d$ general ordinary affine hyperplanes.
They meet each such component in a nonempty finite set in
its good locus, avoiding its intersections with all the
other components. For completeness, this follows by taking
the projective closures: their boundaries and all proper
intersection loci have dimension less than $d$, whereas
$d$ general hyperplanes cut a nonempty zero-dimensional
section of each $d$-dimensional projective component.
The selected points are isolated in the full intersection,
including when larger-dimensional components were present.
Every hyperplane has weighted degree at most $\max_iw_i\leq D$.
We have obtained a square system of $N$ equations of the
required degree, with at least one isolated point for each
original $d$-dimensional component. The square bound gives
\[
       \#\{\text{$d$-dimensional components of }V\}
                         \leq D^N/W.
\]
Summing over $0\leq d\leq N$ proves (D.WB.1).
An open subset of an irreducible scheme is irreducible if
nonempty, so passing to an open cannot increase this number.
Finally every nonempty clopen part contains an irreducible
component, and no irreducible component meets two such parts.
\end{proof}

\begin{corollary}[Proved: the normal-layer product has a quadratic exponent]
\label{D:weighted-normal-layer-budget}
Fix positive constants $A,K$ and an integer $a\geq1$.
Let $m,s,f$ be positive integers, put $r=ms$, and put $q=p^f$
for any fixed prime $p$. Suppose an affine polynomial
presentation has at most $Afms^2$ variables at each weight
$j\in\{1,\ldots,am\}$, and at most $Afms^2$ additional
variables of weight at least one. Suppose the additional
weights are at most $Km$, and every equation has weighted
degree at most $Km$. Enlarging $K$ so that $K\geq a$, its
number of irreducible components is at most
\[
                              q^{C r^2},                 \tag{D.WB.3}
\]
where $C$ depends only on $A,K,a,p$. The same conclusion
holds for any open part and for a partition of it into
nonempty clopen classes.

This statement is conditional on an actual polynomial
presentation with those weights. It does not construct a
presentation of bundles, deformation torsors, or Frobenius
objects from a bound on their tangent spaces.
\end{corollary}

\begin{proof}
Round the fixed constants upward. Pad with unused variables
and equations setting those variables equal to zero until
there are $b=\lceil Afms^2\rceil$ variables at every weight
$1\leq j\leq L=am$ and $b$ additional variables. Padding
does not change the scheme or its components. It does not
violate the degree bound, because every added weight is at
most $Km$. In the denominator of (D.WB.1), conservatively
replace each additional weight by one. If $N=b(L+1)$, the
weighted component bound is at most
\[
        (N+1)\frac{(Km)^{b(L+1)}}{(L!)^b}.              \tag{D.WB.4}
\]
Using $L!\geq(L/e)^L$ yields
\[
 \log\!\left(\frac{(Km)^{b(L+1)}}{(L!)^b}\right)
       \leq bL\log(eK/a)+b\log(Km).
                                                        \tag{D.WB.5}
\]
Both terms are $O(fm^2s^2)$, since $\log m\leq m$ and
$b=O(fms^2)$. Also $\log(N+1)=O(fm^2s^2)$.
Thus (D.WB.4) is at most $\exp(C_0fr^2)$, which is
$q^{(C_0/\log p)r^2}$. This proves (D.WB.3).
In particular, using one common weighted degree $Km$ is
compatible with the bound; it is essential to retain the
product of the variable weights in the denominator.
\end{proof}

\begin{corollary}[Proved: polynomial background coefficients can be separated]
\label{D:weighted-normal-layer-background-budget}
Fix positive constants $A,K,h$ and an integer $a\geq1$.
Let $m,s,f\geq1$, $r=ms$, $q=p^f$, and let $E$ be a
positive integer with $E\leq K(ms)^h$.
Consider a polynomial presentation with the following variables:
\begin{itemize}
\item at most $Afms$ background variables of weight one;
\item at most $Afms^2$ variables at each weight
$Ej$, $1\leq j\leq am$;
\item at most $Afms^2$ further bundle variables of weight $E$.
\end{itemize}
Suppose every equation has weighted degree at most $KEm$,
enlarging $K$ so that $K\geq a$. Then its number of
irreducible components, and the number of nonempty clopen
classes on any open part, is at most $q^{Cr^2}$, with $C$
depending only on $A,K,h,a,p$.

In particular, this numerical bound permits polynomially
growing degrees in a smaller block of background variables,
provided the actual mixed polynomials satisfy the displayed
weighted-degree condition. That condition is not implied
by separate tangent-space dimensions.
\end{corollary}

\begin{proof}
As before, round the constants upward and pad the three
blocks by zero variables. Put
$b=\lceil Afms^2\rceil$, $c=\lceil Afms\rceil$, and $L=am$.
There are $N=b(L+1)+c$ variables. Lemma
\ref{D:weighted-affine-component-bound} bounds the number
of components by
\[
 (N+1)\frac{(KEm)^N}{E^{b(L+1)}(L!)^b}
   =(N+1)\frac{(Km)^{b(L+1)}}{(L!)^b}(KEm)^c.
                                                        \tag{D.WB.6}
\]
The first fraction has logarithm $O(fm^2s^2)$ by the
calculation in (D.WB.5). The extra logarithm is
\[
 c\log(KEm)
       =O\bigl(fms(1+\log(ms))\bigr)
       =O(fm^2s^2),                                    \tag{D.WB.7}
\]
using $E\leq K(ms)^h$ and $1+\log x\leq 2x$ for $x\geq1$.
Also $\log(N+1)=O(fm^2s^2)$. Conversion from
$\exp(O(fr^2))$ to $q^{Cr^2}$ gives the result.
\end{proof}
